% =============================================================================
% PRLab Tech Report Template — main.tex
% Compile with pdfLaTeX (e.g. on Overleaf, or locally via `pdflatex main`).
% Class is a verbatim fork of the Tsinghua College of AI / ByteDance Seed
% tech-report template, with PRLab purple #670095 in place of Seed blue.
% =============================================================================
\documentclass[]{prlab}

\usepackage[toc,page,header]{appendix}
\microtypesetup{expansion=false}
\usepackage{graphicx}
\usepackage{subcaption}
\usepackage{booktabs}
\usepackage{hyperref}
\usepackage{bm}
\usepackage{amsmath}
\usepackage{amssymb}
\usepackage{mathtools}
\usepackage{amsthm}
\usepackage{xspace}
\usepackage{wrapfig}
\usepackage{multirow}
\usepackage{tabularx}
\usepackage{colortbl}
\usepackage{pifont}        % \ding for the \cmark / \xmark markers used in sec/
\normalem  % the class loads ulem (\emph defaults to underline) — restore italic

% Tick / cross markers used in the example sections.
\newcommand{\cmark}{\textcolor{green!60!black}{\ding{51}}}
\newcommand{\xmark}{\textcolor{red!75!black}{\ding{55}}}


\title{PosterCopilot: Toward Layout Reasoning and Controllable Editing for Professional Graphic Design}

\author[1,\star]{First Author}
\author[1,\star]{Second Author}
\author[1]{Third Author}
\author[2,\dagger]{Fourth Author}

\affiliation[1]{Pattern Recognition Laboratory, Nanjing University}
\affiliation[2]{Collaborating Institution}

\contribution[\star]{Equal Contribution}
\contribution[\dagger]{Corresponding Author}

\date{\today}
\checkdata[Code]{\url{https://github.com/your-org/your-project}}
\checkdata[Project Page]{\url{https://your-project.example.com}}
\correspondence{\email{correspondence@nju.edu.cn}}

\abstract{%
Graphic design forms the cornerstone of modern visual communication, serving as a vital medium for promoting cultural and commercial events. Recent advances have explored automating this process using Large Multimodal Models (LMMs), yet existing methods often produce geometrically inaccurate layouts and lack the iterative, layer-specific editing required in professional workflows. To address these limitations, we present PosterCopilot, a framework that advances layout reasoning and controllable editing for professional graphic design. Specifically, we introduce a progressive three-stage training strategy that equips LMMs with geometric understanding and aesthetic reasoning for layout design, consisting of Perturbed Supervised Fine-Tuning, Reinforcement Learning for Visual-Reality Alignment, and Reinforcement Learning from Aesthetic Feedback. Furthermore, we develop a complete workflow that couples the trained LMM-based design model with generative models, enabling layer-controllable, iterative editing for precise element refinement while maintaining global visual consistency. Extensive experiments demonstrate that PosterCopilot achieves geometrically accurate and aesthetically superior layouts, offering fine-grained, layer-wise controllability for professional iterative design. 
}


\begin{document}
\maketitle

\section{Introduction}
\label{sec:intro}

This template provides a clean starting point for tech reports out of the
Pattern Recognition Laboratory at Nanjing University. Section headings are
rendered in PRLab purple (\texttt{\#3B1A6B}) using IBM Plex Sans, while
the body uses IBM Plex Serif. Citations such as \citep{vaswani2017attention}
appear as numeric references in the same purple, giving the document a
consistent visual rhythm.

\paragraph{How to use this template.}
Edit \texttt{main.tex} to set your title, authors, affiliations, abstract,
and metadata. Then add or rewrite the section stubs under \texttt{sec/}.
Drop your figures into \texttt{fig/} and your BibTeX entries into
\texttt{references.bib}. Compile with pdfLaTeX:

\begin{center}
\texttt{pdflatex main \&\& bibtex main \&\& pdflatex main \&\& pdflatex main}
\end{center}

\paragraph{What the template gives you.}
A first-page header with PRLab and Nanjing University wordmarks, a framed
abstract block, three-level section headings in PRLab purple, sans-serif
captions with colored labels, numeric citations in PRLab purple, and a
matching appendix style. The wordmarks in \texttt{prlab.cls} are TikZ
placeholders --- swap in real logo assets via \texttt{\textbackslash
includegraphics} when they are available.

\section{Method}
\label{sec:method}

We describe the method here. Mathematics is set with \texttt{newpxmath}, which
pairs cleanly with IBM Plex Serif body text. For example, given an input
$\mathbf{x} \in \mathbb{R}^d$, we compute
\begin{equation}
  \mathbf{y} \;=\; \sigma(\mathbf{W}\mathbf{x} + \mathbf{b}),
  \label{eq:linear}
\end{equation}
where $\sigma$ is the activation. Refer to \cref{eq:linear} via
\texttt{cleveref}, which is preloaded.

\subsection{Architecture}
\label{sec:arch}

Subsection headings are slightly smaller than section headings but use the
same PRLab purple. Inline emphasis still works as expected:
\textbf{bold text} renders in sans-serif (matching the Seed and Meta
conventions), while \emph{italic} stays in the serif body face.

\subsection{Training Details}
\label{sec:train}

A representative table looks like this:

\begin{table}[h]
  \centering
  \begin{tabular}{lcc}
    \toprule
    \textbf{Model} & \textbf{Params} & \textbf{Accuracy} \\
    \midrule
    Baseline & 12M  & 71.3 \\
    Ours     & 14M  & \textbf{78.9} \\
    Ours-XL  & 48M  & \textbf{82.4} \\
    \bottomrule
  \end{tabular}
  \caption{A representative comparison table.\label{tab:main}}
\end{table}

The caption above demonstrates the PRLab caption style: a sans-serif purple
label (``Table~1'') followed by the caption text in regular small body.

\section{Experiments}
\label{sec:experiments}

\subsection{Setup}

We follow standard protocols. References look like
\citep{he2016deep,dosovitskiy2020image}, where the bracketed numbers are
rendered in PRLab purple.

\subsection{Main Results}

\subsubsection{Subsubsection Example}

Three-level section nesting works out of the box. This is the deepest level
that receives the purple sans-serif treatment.

\paragraph{Inline paragraph headings.}
Paragraph headings are italic and stay inline with the running text, in line
with the Seed template convention.

\begin{itemize}
  \item First bulleted observation.
  \item Second bulleted observation with a \cmark{} marker.
  \item Third bulleted observation with an \xmark{} marker.
\end{itemize}

\section{Conclusion}
\label{sec:conclusion}

This section is where you summarize contributions and outline directions for
future work. Keep it brief.


\clearpage
\bibliographystyle{plainnat}
\bibliography{references}

\clearpage
\beginappendix
\section{Additional Details}
\label{app:details}

Use the appendix for derivations, extra experiments, prompts, and any other
material that would interrupt the main narrative.

\subsection{Notation}

A subsection in the appendix renders identically to one in the main body,
following the PRLab purple sans-serif convention.


\end{document}
